% GENERATED FILE. Edit canonical lessons, then run npm run generate:books.
% canonical-source-hash: fnv1a64:854c354069437a8d
% canonical-lessons: LA-C326-adulescentia, LA-C326-iuventus, LA-C326-senectus, LA-C326-orbus, LA-C326-pupillus

\chapter{Aetātēs}
\label{ch:la-aetates}
\hlchaptermodality{326}

\begin{chapteropening}
\textbf{By the end of this chapter:} \emph{I can talk about the ages of life: young adulthood, youth, old age, orphaned and a ward.}

Complete the last lesson of chapter 326: I can talk about the ages of life: young adulthood, youth, old age, orphaned and a ward.
\end{chapteropening}

\section[adulēscentia, adulēscentiae]{adulēscentia, adulēscentiae — young adulthood, youth}

\label{lesson:LA-C326-adulescentia}



\begin{quote}
Before the new one: say the Latin for an extra family name, then the Latin for a guardian.
\end{quote}

\subsection*{You'll want to know: adulēscentia, adulēscentiae}
\textbf{adulēscentia, adulēscentiae} — \textquotedblleft{}young adulthood, youth\textquotedblright{}.

\begin{quote}
In adulēscentiā multa et stulta fēcī.
\end{quote}

\begin{quote}
\emph{In my youth I did many foolish things.}
\end{quote}

Feminine, from \emph{adulēscēns} (taught): roughly from fifteen to thirty. English \emph{adolescence}.

\begin{cousinweb}
\textbf{adulēscentia} is built on \emph{adulēscēns}, the participle of \emph{adulēscō}, \textquotedblleft{}to grow up\textquotedblright{} (Lewis \& Short). English \emph{adolescence} comes from it.
\end{cousinweb}

\begin{grammarlens}[title={What you've built}]
One more word for telling, explaining and describing at B1.
\end{grammarlens}

\subsection*{Your turn}
\begin{itemize}
\raggedright
  \item \emph{Say it:} \emph{adulēscentia, adulēscentiae}
  \item \emph{Say it:} \emph{adulēscentia, adulēscentiae}, once more
  \item \emph{Say it:} a sentence of your own with \emph{adulēscentia, adulēscentiae}
  \item \emph{Recall:} say the Latin for an extra family name, then the Latin for a guardian, then say \emph{adulēscentia, adulēscentiae} again
\end{itemize}

\begin{tcolorbox}[breakable,colback=teal!4,colframe=teal!35!black,title={Before you move on}]
What does adulēscentia, adulēscentiae mean? (\textquotedblleft{}Young adulthood, youth\textquotedblright{}.) Say it once more.
\end{tcolorbox}
\section[iuventūs, iuventūtis]{iuventūs, iuventūtis — youth; young people}

\label{lesson:LA-C326-iuventus}



\begin{quote}
Before the new one: say the Latin for a guardian, then the Latin for young adulthood.
\end{quote}

\subsection*{You'll want to know: iuventūs, iuventūtis}
\textbf{iuventūs, iuventūtis} — \textquotedblleft{}youth; young people\textquotedblright{}.

\begin{quote}
Tōta iuventūs urbis in campō lūdēbat.
\end{quote}

\begin{quote}
\emph{All the young people of the city were playing on the field.}
\end{quote}

Feminine, 3rd declension, from \emph{iuvenis}, 'young'. Both the time of youth and the young people themselves.

\begin{cousinweb}
\textbf{iuventūs} is built on \emph{iuvenis}, \textquotedblleft{}young\textquotedblright{} (EtymDB), which is inherited and a cousin of English \emph{young}. English \emph{juvenile} comes from \emph{iuvenis}.
\end{cousinweb}

\begin{grammarlens}[title={What you've built}]
One more word for telling, explaining and describing at B1.
\end{grammarlens}

\subsection*{Your turn}
\begin{itemize}
\raggedright
  \item \emph{Say it:} \emph{iuventūs, iuventūtis}
  \item \emph{Say it:} \emph{iuventūs, iuventūtis}, once more
  \item \emph{Say it:} a sentence of your own with \emph{iuventūs, iuventūtis}
  \item \emph{Recall:} say the Latin for a guardian, then the Latin for young adulthood, then say \emph{iuventūs, iuventūtis} again
\end{itemize}

\begin{tcolorbox}[breakable,colback=teal!4,colframe=teal!35!black,title={Before you move on}]
What does iuventūs, iuventūtis mean? (\textquotedblleft{}Youth; young people\textquotedblright{}.) Say it once more.
\end{tcolorbox}
\section[senectūs, senectūtis]{senectūs, senectūtis — old age}

\label{lesson:LA-C326-senectus}



\begin{quote}
Before the new one: say the Latin for young adulthood, then the Latin for youth.
\end{quote}

\subsection*{You'll want to know: senectūs, senectūtis}
\textbf{senectūs, senectūtis} — \textquotedblleft{}old age\textquotedblright{}.

\begin{quote}
Avus meus in senectūte in vīllā habitābat.
\end{quote}

\begin{quote}
\emph{In his old age my grandfather lived at the villa.}
\end{quote}

Feminine, 3rd declension, from \emph{senex} (taught). Cicero wrote a whole dialogue in its praise.

\begin{cousinweb}
\textbf{senectūs} is built on \emph{senex}, \textquotedblleft{}old man\textquotedblright{} (Lewis \& Short). English \emph{senior} and \emph{senile} come from the same family.
\end{cousinweb}

\begin{grammarlens}[title={What you've built}]
One more word for telling, explaining and describing at B1.
\end{grammarlens}

\subsection*{Your turn}
\begin{itemize}
\raggedright
  \item \emph{Say it:} \emph{senectūs, senectūtis}
  \item \emph{Say it:} \emph{senectūs, senectūtis}, once more
  \item \emph{Say it:} a sentence of your own with \emph{senectūs, senectūtis}
  \item \emph{Recall:} say the Latin for young adulthood, then the Latin for youth, then say \emph{senectūs, senectūtis} again
\end{itemize}

\begin{tcolorbox}[breakable,colback=teal!4,colframe=teal!35!black,title={Before you move on}]
What does senectūs, senectūtis mean? (\textquotedblleft{}Old age\textquotedblright{}.) Say it once more.
\end{tcolorbox}
\section[orbus, orba, orbum]{orbus, orba, orbum — orphaned; childless, bereft}

\label{lesson:LA-C326-orbus}



\begin{quote}
Before the new one: say the Latin for youth, then the Latin for old age.
\end{quote}

\subsection*{You'll want to know: orbus, orba, orbum}
\textbf{orbus, orba, orbum} — \textquotedblleft{}orphaned; childless, bereft\textquotedblright{}.

\begin{quote}
Puella orba ab avunculō ēducāta est.
\end{quote}

\begin{quote}
\emph{The orphaned girl was brought up by her uncle.}
\end{quote}

Bereft of parents or of children; with the ablative, bereft of anything.

\begin{cousinweb}
\textbf{orbus} is an inherited word that Lewis \& Short compare with Greek \emph{orphanos}, \textquotedblleft{}orphan\textquotedblright{}, the source of English \emph{orphan}.
\end{cousinweb}

\begin{grammarlens}[title={What you've built}]
One more word for telling, explaining and describing at B1.
\end{grammarlens}

\subsection*{Your turn}
\begin{itemize}
\raggedright
  \item \emph{Say it:} \emph{orbus, orba, orbum}
  \item \emph{Say it:} \emph{orbus, orba, orbum}, once more
  \item \emph{Say it:} a sentence of your own with \emph{orbus, orba, orbum}
  \item \emph{Recall:} say the Latin for youth, then the Latin for old age, then say \emph{orbus, orba, orbum} again
\end{itemize}

\begin{tcolorbox}[breakable,colback=teal!4,colframe=teal!35!black,title={Before you move on}]
What does orbus, orba, orbum mean? (\textquotedblleft{}Orphaned; childless, bereft\textquotedblright{}.) Say it once more.
\end{tcolorbox}
\section[pūpillus, pūpillī]{pūpillus, pūpillī — a ward, an orphan boy}

\label{lesson:LA-C326-pupillus}



\begin{quote}
Before the new one: say the Latin for old age, then the Latin for orphaned.
\end{quote}

\subsection*{You'll want to know: pūpillus, pūpillī}
\textbf{pūpillus, pūpillī} — \textquotedblleft{}a ward, an orphan boy\textquotedblright{}.

\begin{quote}
Tūtor pecūniam pūpillī fidēliter servābat.
\end{quote}

\begin{quote}
\emph{The guardian faithfully kept the ward's money safe.}
\end{quote}

Masculine. The girl is a \emph{pūpilla}, also 'the pupil of the eye'. English \emph{pupil}.

\begin{cousinweb}
\textbf{pūpillus} is a diminutive of \emph{pūpulus}, \textquotedblleft{}little child\textquotedblright{} (Lewis \& Short): a little child in the care of another. English \emph{pupil} comes from it through French.
\end{cousinweb}

\begin{grammarlens}[title={What you've built}]
One more word for telling, explaining and describing at B1.
\end{grammarlens}

\subsection*{Your turn}
\begin{itemize}
\raggedright
  \item \emph{Say it:} \emph{pūpillus, pūpillī}
  \item \emph{Say it:} \emph{pūpillus, pūpillī}, once more
  \item \emph{Say it:} a sentence of your own with \emph{pūpillus, pūpillī}
  \item \emph{Recall:} say the Latin for old age, then the Latin for orphaned, then say \emph{pūpillus, pūpillī} again
\end{itemize}

\begin{tcolorbox}[breakable,colback=teal!4,colframe=teal!35!black,title={Before you move on}]
What does pūpillus, pūpillī mean? (\textquotedblleft{}A ward, an orphan boy\textquotedblright{}.) Say it once more.
\end{tcolorbox}
